% Chapter 27: Mediation Analysis
% A First Course in Causal Inference - Peng Ding

\chapter{中介分析：自然直接效应与自然间接效应}
\label{ch:mediation}

\section{本章导论：一个古老的追问}

当我们说``吸烟导致肺癌''时，这个说法究竟意味着什么？

一方面，烟草中的有害物质可能直接作用于细胞，引发癌变；另一方面，吸烟可能通过影响其他生理指标——比如血压、免疫力或慢性炎症——进而间等地增加患癌风险。这两种机制，本质上完全不同：一种是直接杀伤，另一种是通过一连串中间步骤传导的\textbf{间接效应}。

理解因果效应的这种\textbf{传导机制}（transmission mechanism），是科学研究的核心关切之一。生物医学研究者想知道某种药物是通过什么生化通路发挥作用的；社会学家想知道贫困是通过什么社会机制影响健康结果的；政策制定者想知道一项教育改革是直接改善了学生能力，还是通过改变学校资源配置间接起效。\textbf{中介分析（mediation analysis）}正是回答这类问题的统计框架。

本章的组织如下。我们首先通过具体实例说明中介分析的应用场景；然后引入\textbf{嵌套潜在结果}（nested potential outcomes）的数学框架；接着建立识别中介效应的\textbf{中介公式}（mediation formula）；最后在线性模型下推导出经典的\textbf{Baron--Kenny 公式}，并讨论其敏感性分析。

\section{从总效应到直接效应与间接效应}

在中介分析的框架中，我们有四个核心变量：处理 $Z$（treatment）、结果 $Y$（outcome）、中介 $M$（mediator）以及背景协变量 $X$（covariates）。它们之间的因果关系可以直观地用有向无环图（DAG）表示（\cref{fg:mediation-dag}）。

\begin{figure}[ht]
\centering
\begin{xy}
\xymatrix{
&& X \ar[dll] \ar[d] \ar[drr] \\
Z \ar[rr] \ar@/_1.5pc/[rrrr] & & M\ar[rr] & & Y}
\end{xy}
\caption{中介分析的因果有向无环图：$Z$ 通过 $M$ 影响 $Y$ 的路径，以及 $Z$ 直接影响 $Y$ 的路径。}
\label{fg:mediation-dag}
\end{figure}

\fboxsep=6pt\fboxrule=1pt
\noindent\fbox{
\begin{minipage}{0.97\textwidth}
\textbf{核心问题}：总的因果效应 $\tau$ 中，有多少是通过中介变量 $M$ 传导的（\textbf{间接效应}），有多少是不经过 $M$ 的（\textbf{直接效应}）？
\end{minipage}
}

以下三个真实研究场景可以帮助我们理解中介分析的应用。

\begin{Example}[遗传变异与肺癌：通过吸烟的中介]
\Vanderweele{}（2012）研究了染色体 15q25.1 上的遗传变异如何影响肺癌风险。\footnote{详见 \cite{Vanderweele::2012genetic}。} 该研究的核心问题是：遗传变异是通过影响吸烟行为（吸烟量）进而增加肺癌风险的，还是通过其他独立通路（如直接影响肺部细胞）直接致癌的？这就是一个典型的中介分析问题：遗传变异是处理 $Z$，吸烟是中介 $M$，肺癌是结果 $Y$。
\end{Example}

\begin{Example}[邻里贫困与青少年物质使用：学校与同伴环境的中介]
\Rudolf{} 等人（2018）研究了邻里贫困对青少年物质使用的影响，并探讨了这一影响是通过学校和同伴环境传导的。\footnote{详见 \cite{rudolph2018causal}。} 他们使用了美国国家共病调查（National Comorbidity Survey）的数据，将邻里贫困指数作为处理，青少年物质使用（如饮酒、吸毒）作为结果，而学校和同伴环境因素作为中介。
\end{Example}

\begin{Example}[就业培训对心理健康：工作搜寻效能的中介]
JOBS II 是一个关于就业培训对失业工人影响的现场实验。\footnote{我们在 \cref{sec:cace-application} 章中曾使用过这个数据集。} 该培训项目的目的不仅是提高再就业率，还希望改善求职者的心理健康。研究者感兴趣的是：培训对心理健康的效应中，有多少是通过提高工作搜寻效能这一中介传导的，又有多少是通过其他途径（如社会支持网络）直接产生的。\footnote{我们将在 \cref{sec:bk-jobs} 节中回到这个例子，展示 Baron--Kenny 方法的具体应用。}
\end{Example}

\section{嵌套潜在结果}

\subsection{潜在结果的两层结构}

为了定义直接效应与间接效应，我们需要引入\textbf{嵌套潜在结果}（nested potential outcomes）这一概念。\Citet{Robins::1992} 和 \Citet{Pearl::2001} 是这一理论的奠基人。

首先，考虑对处理 $z$ 进行干预，定义对应的潜在中介和潜在结果：
\[
\{ M(z), \; Y(z) : z = 0, 1 \}.
\]
然后，考虑同时对 $z$ 和 $m$ 进行干预，定义：
\[
\{ Y(z, m) : z = 0, 1, \; m \in \mathcal{M} \},
\]
其中 $\mathcal{M}$ 是中介所有可能的取值范围。

关键的一步是考虑\textbf{嵌套干预}（nested intervention）：将 $z$ 设定为某个值，但将 $m$ 设定为在另一个处理水平下的潜在中介值。记 $M_{z'} \equiv M(z')$，\Citet{Robins::1992} 和 \Citet{Pearl::2001} 定义了嵌套潜在结果：
\[
\{ Y(z, M_{z'}) : z = 0, 1, \; z' = 0, 1 \}.
\]

其中 $Y(z, M_{z'})$ 表示：如果治疗被设定为 $z$，而中介被设定为在治疗 $z'$ 下的潜在值 $M_{z'}$，则结果会是什么。注意 $z$ 和 $z'$ 可以不同——这正是其``嵌套''之处。

在二值处理的设定下，我们有四个嵌套潜在结果：
\[
\{ Y(1, M_1), \; Y(1, M_0), \; Y(0, M_1), \; Y(0, M_0) \}.
\]

其中 $Y(1, M_1)$ 表示：如果治疗设为 $z=1$ 且中介设为在 $z=1$ 下的自然值，则结果如何。类似地，$Y(0, M_0)$ 表示治疗设为 $z=0$ 且中介设为在 $z=0$ 下的自然值的情形。

Composition 假设确保了嵌套潜在结果与标准潜在结果之间的一致性。

\begin{assumption}[Composition 假设]\label{assume:composition}
$Y(z, M_z) = Y(z)$，对 $z = 0, 1$ 成立。
\end{assumption}

Composition 假设是不可证明的，它本质上是一个假设。但在实践中，我们完全可以将 $Y(1)$ 定义为 $Y(1, M_1)$，将 $Y(0)$ 定义为 $Y(0, M_0)$，这样该假设就自动成立了。

现在，我们可以正式定义三种效应。

\begin{Definition}[总效应、自然直接效应与自然间接效应]\label{def:mediation}
总效应（total effect）定义为
\[
\tau = \mathbb{E}\{ Y(1) - Y(0) \}.
\]

自然直接效应（natural direct effect, NDE）定义为
\[
\NDE = \mathbb{E}\{ Y(1, M_0) - Y(0, M_0) \}.
\]

自然间接效应（natural indirect effect, NIE）定义为
\[
\NIE = \mathbb{E}\{ Y(1, M_1) - Y(1, M_0) \}.
\]
\end{Definition}

这三个效应的含义如下：

\begin{itemize}
  \item \textbf{总效应} $\tau$：标准的平均因果效应，即治疗从 $0$ 变到 $1$ 对结果的平均影响。
  \item \textbf{NDE} $\mathbb{E}\{Y(1, M_0) - Y(0, M_0)\}$：如果将治疗从 $0$ 变到 $1$，但把中介固定在无治疗时的自然水平 $M_0$，结果的平均变化。这个效应衡量的是\textbf{不经过 $M$ 的直接通路}的贡献。
  \item \textbf{NIE} $\mathbb{E}\{Y(1, M_1) - Y(1, M_0)\}$：如果保持治疗 $z=1$ 不变，但把中介从无治疗时的自然水平 $M_0$ 变到治疗时的自然水平 $M_1$，结果的平均变化。这个效应衡量的是\textbf{通过中介 $M$ 传导}的间接贡献。
\end{itemize}

在 Composition 假设下，总效应可以分解为自然直接效应与自然间接效应之和。

\begin{Proposition}[效应的分解]\label{prop:mediation-decomposition}
在 \cref{def:mediation} 和 \cref{assume:composition} 下，有
\[
\tau = \NDE + \NIE.
\]
\end{Proposition}

\fboxsep=5pt\fboxrule=1pt
\noindent\fbox{
\begin{minipage}{0.97\textwidth}
\textbf{为什么 NIE 可以衡量间接效应？}直观上，$Y(1, M_1) - Y(1, M_0)$ 表示：将治疗固定在 $z=1$，但让中介从``如果没有治疗它会是什么''变为``如果有治疗它会是什么''，结果会如何变化。这恰恰描述了``治疗通过改变中介进而影响结果''这条通路上的效应。
\end{minipage}
}

然而，$Y(1, M_0)$ 这个量本身有一个深刻的概念困难：它是一个\textbf{跨世界counterfactual}（cross-world counterfactual）。$z=1$ 和 $m=M_0$ 这两个干预不可能在同一个实验中同时发生——同一个单元不可能既接受治疗 $z=1$，又同时处于无治疗条件下的中介水平。要理解 $Y(1, M_0)$，我们必须借助平行世界的想象（\cref{fg:crossworld}）。

\begin{figure}[ht]
\centering
[width = 0.7\textwidth]{figures/crossworld.pdf}
\caption{跨世界 counterfactual $Y(1, M_0)$ 和 $Y(0, M_1)$ 的示意图。同一单元在两个平行世界中分别接受不同的治疗，从而获得不同的中介值。}
\label{fg:crossworld}
\end{figure}

这种跨世界的性质引发了一个哲学上的争议。\Citet{frangakis2002principal} 将 $Y(1, M_0)$ 和 $Y(0, M_1)$ 称为\textbf{先验 counterfactual}（a priori counterfactuals），因为我们无法在任何物理实验中观察到它们。按照 \Citet{popper1963conjectures} 的可证伪性标准，如果一个陈述无法被任何物理实验或观测所证伪，那它就不是科学的，而是形而上学的。因此，严格的 Popperian 统计学家可能会认为中介分析是形而上学。\footnote{这一争议的详细讨论见 \cite{dawid2000causal} 和 \cite{robins2010treatment}。}

不过，\Citet{dawid2000causal} 对整个潜在结果框架都提出了类似的批评。但他的批评是过度的：即使我们无法同时观察到 $Y(1)$ 和 $Y(0)$，我们仍然可以通过边缘分布的信息对它们做出可证伪的陈述。$Y(1, M_0)$ 的困难在于它是真正无法在任何实验中触及的——这与 $\{Y(1), Y(0)\}$ 的性质根本不同。

\section{中介公式}

\subsection{四个关键假设}

中介公式的识别依赖于四个假设。前三个假设本质上是在说：治疗和中介在给定协变量 $X$ 的条件下都是被随机化的。

\begin{assumption}[无治疗—结果混淆]\label{assume:mediation-no-z-y-confounding}
给定 $X$，治疗与潜在结果独立：
\[
Z \ind Y(z, m) \mid X, \quad \forall z, m.
\]
\end{assumption}

\begin{assumption}[无中介—结果混淆]\label{assume:mediation-no-m-y-confounding}
给定 $(X, Z)$，中介与潜在结果独立：
\[
M \ind Y(z, m) \mid (X, Z), \quad \forall z, m.
\]
\end{assumption}

\begin{assumption}[无治疗—中介混淆]\label{assume:mediation-no-z-m-confounding}
给定 $X$，治疗与潜在中介独立：
\[
Z \ind M(z) \mid X, \quad \forall z.
\]
\end{assumption}

这三个假设共同构成了\textbf{序列可忽略性}（sequential ignorability）：在给定 $X$ 的条件下，治疗分配是可忽略的；进一步在给定 $(X, Z)$ 的条件下，中介分配也是可忽略的。序列可忽略性等价于在给定 $X$ 的条件下 $(Z, M)$ 是联合随机化的：
\begin{equation}
(Z, M) \ind Y(z, m) \mid X, \quad \forall z, m.
\label{eq:joint-randomization}
\end{equation}
\cref{eq:joint-randomization} 与 \cref{assume:mediation-no-z-y-confounding,assume:mediation-no-m-y-confounding} 的等价性留作练习。\footnote{详见 \cref{ex:sr-jr-equivalence} 题。}

最后一个假设是跨世界独立性。

\begin{assumption}[跨世界独立性]\label{assume:mediation-crossworld}
给定 $X$，潜在结果与潜在中介独立：
\[
Y(z, m) \ind M(z') \mid X, \quad \forall z, z', m.
\]
\end{assumption}

前三个假设至少在随机化实验中是可以保证的，但 \cref{assume:mediation-crossworld} 更为强——没有任何物理实验可以确保它，因为当 $z \neq z'$ 时，我们永远无法同时观察到 $Y(z, m)$ 和 $M(z')$。这正是 $Y(1, M_0)$ 等 cross-world counterfactual 的根本困难所在。

\begin{Example}[验证假设的数据生成过程]\label{eg:mediation-assumptions}
给定 $X$，我们按如下方式生成数据：
\begin{align*}
Z &= 1\{ g_Z(X, \varepsilon_Z) \geq 0 \}, \\
M(z) &= 1\{ g_M(X, z, \varepsilon_M) \geq 0 \}, \quad z = 0, 1, \\
Y(z, m) &= g_Y(X, z, m, \varepsilon_Y), \quad z, m = 0, 1,
\end{align*}
其中 $\varepsilon_Z, \varepsilon_M, \varepsilon_Y$ 相互独立。\footnote{详见 \cite{robins2010treatment}。} 可以验证，在上述数据生成过程下，\cref{assume:mediation-no-z-y-confounding,assume:mediation-no-m-y-confounding,assume:mediation-no-z-m-confounding,assume:mediation-crossworld} 均成立。如果允许 $\varepsilon_M$ 和 $\varepsilon_Y$ 依赖于处理水平（即改为 $\varepsilon_M(z)$ 和 $\varepsilon_Y(z, m)$），则这些假设可能被违反。
\end{Example}

\subsection{识别公式}

\Citet{Pearl::2001} 证明了以下关键结果。

\begin{Theorem}[中介公式]\label{eq:mediation-formula}
在 \cref{assume:mediation-no-z-y-confounding,assume:mediation-no-m-y-confounding,assume:mediation-no-z-m-confounding,assume:mediation-crossworld} 下，有
\[
\mathbb{E}\{ Y(z, M_{z'}) \mid X = x \}
= \sum_m \mathbb{E}(Y \mid Z = z, M = m, X = x) \; \Pr(M = m \mid Z = z', X = x).
\]
\end{Theorem}

\fboxsep=5pt\fboxrule=1pt
\noindent\fbox{
\begin{minipage}{0.97\textwidth}
\textbf{定理的含义}：我们无法直接观察跨世界 counterfactual $Y(z, M_{z'})$，但 \cref{assume:mediation-no-z-y-confounding,assume:mediation-no-m-y-confounding,assume:mediation-no-z-m-confounding,assume:mediation-crossworld} 这四个假设足够将这个不可观测的量表达为可观测条件分布的函数。具体来说，$Y(z, M_{z'})$ 的期望等于：结果的条件期望（给定 $Z=z, M=m, X=x$）对中介的分布（给定 $Z=z', X=x$）的加权和。
\end{minipage}
}

\fboxsep=5pt\fboxrule=1pt
\noindent\fbox{
\begin{minipage}{0.97\textwidth}
\textbf{为什么四个假设缺一不可？}定理的证明（见附录 \cref{sec:proof-mediation}）清楚地展示了这四个假设各自的角色：\cref{assume:mediation-no-z-m-confounding} 保证了 $\Pr(M = m \mid Z = z', X = x) = \Pr(M_{z'} = m \mid X = x)$，\cref{assume:mediation-crossworld} 允许我们交换 $Y(z, m)$ 与 $M_{z'}$ 的条件，\cref{assume:mediation-no-z-y-confounding,assume:mediation-no-m-y-confounding} 则保证了结果的识别。
\end{minipage}
}

对 $X$ 求期望，得到无条件版本：
\[
\mathbb{E}\{ Y(z, M_{z'}) \}
= \sum_x \sum_m \mathbb{E}(Y \mid Z = z, M = m, X = x) \; \Pr(M = m \mid Z = z', X = x) \; \Pr(X = x).
\]

对于连续型 $M$ 和 $X$，求和替换为积分。

由 \cref{prop:mediation-decomposition} 和 \cref{eq:mediation-formula}，我们可以得到 NDE 和 NIE 的识别公式。

\begin{Corollary}[NDE 和 NIE 的识别]\label{cor:nde-nie}
在 \cref{assume:mediation-no-z-y-confounding,assume:mediation-no-m-y-confounding,assume:mediation-no-z-m-confounding,assume:mediation-crossworld} 下，条件 NDE 和 NIE 由下式识别：
\begin{align}
\NDE(x)
&= \sum_m \big[ \mathbb{E}(Y \mid Z = 1, M = m, X = x) - \mathbb{E}(Y \mid Z = 0, M = m, X = x) \big] \Pr(M = m \mid Z = 0, X = x), \label{eq:nde-formula} \\
\NIE(x)
&= \sum_m \mathbb{E}(Y \mid Z = 1, M = m, X = x) \; \big[ \Pr(M = m \mid Z = 1, X = x) - \Pr(M = m \mid Z = 0, X = x) \big]. \label{eq:nie-formula}
\end{align}
\end{Corollary}

\fboxsep=5pt\fboxrule=1pt
\noindent\fbox{
\begin{minipage}{0.97\textwidth}
\textbf{理解 \cref{eq:nde-formula} 的直觉}：$\NDE(x)$ 衡量的是将处理从 $0$ 变到 $1$（但保持中介在无处理时的自然分布）的效应。在每个中介值 $m$ 下，处理从 $0$ 变到 $1$ 的条件效应是 $\mathbb{E}(Y \mid Z=1, M=m, X=x) - \mathbb{E}(Y \mid Z=0, M=m, X=x)$；由于中介被固定在 $Z=0$ 条件下的分布，这些条件效应需要按 $\Pr(M=m \mid Z=0, X=x)$ 加权求和。
\end{minipage}
}

对于二值中介 $M \in \{0, 1\}$，\cref{cor:nde-nie} 中的 NIE 公式可以化简为一种优美的\textbf{乘积形式}。

\begin{Corollary}[二值中介的 NIE 乘积公式]\label{cor:nie-binary}
在 \cref{assume:mediation-no-z-y-confounding,assume:mediation-no-m-y-confounding,assume:mediation-no-z-m-confounding,assume:mediation-crossworld} 下，对二值中介 $M$，有
\[
\NIE(x) = \tau_{Z \to M}(x) \cdot \tau_{M \to Y}(1, x),
\]
其中
\[
\tau_{Z \to M}(x) = \Pr(M = 1 \mid Z = 1, X = x) - \Pr(M = 1 \mid Z = 0, X = x)
\]
是 $Z$ 对 $M$ 的条件平均因果效应，
\[
\tau_{M \to Y}(z, x) = \mathbb{E}(Y \mid Z = z, M = 1, X = x) - \mathbb{E}(Y \mid Z = z, M = 0, X = x)
\]
是 $M$ 对 $Y$ 的条件平均因果效应。
\end{Corollary}

\fboxsep=5pt\fboxrule=1pt
\noindent\fbox{
\begin{minipage}{0.97\textwidth}
\textbf{为什么叫乘积公式？}直观上，间接效应是这样传导的：首先，处理 $Z$ 变化导致中介 $M$ 变化（效应 $\tau_{Z \to M}$）；然后，中介 $M$ 变化导致结果 $Y$ 变化（效应 $\tau_{M \to Y}$）。这两段效应的乘积，就是总间接效应。这与我们的直觉完全一致。
\end{minipage}
}

\Cref{cor:nie-binary} 的证明留作练习。\footnote{详见 \cref{ex:nie-binary-m} 题。}

\section{线性模型下的 Baron--Kenny 方法}

\cref{eq:mediation-formula} 给出了非参数识别公式，它允许我们在各种不同模型下推导出中介效应的具体形式。本节介绍经典而又广泛使用的\textbf{Baron--Kenny 方法}（\citeyear{baron1986moderator}），该方法假设线性模型。

\subsection{模型设定}

Baron--Kenny 方法假设中介和结果都满足线性模型（\cref{fg:mediation-dag-bk}）。

\begin{figure}[ht]
\centering
$$
\begin{xy}
\xymatrix{
&  & X \ar[ddll] \ar[dd]^{\beta_2} \ar[ddrr]^{\theta_4} \\
&&\\
Z \ar[rr]^{\beta_1} \ar@/_2.8pc/[rrrr]_{\theta_1} & & M\ar[rr]^{\theta_2} & & Y}
\end{xy}
$$
\caption{Baron--Kenny 方法的路径图：$\beta_1 \theta_2$ 为间接效应，$\theta_1$ 为直接效应。}
\label{fg:mediation-dag-bk}
\end{figure}

\begin{assumption}[Baron--Kenny 线性模型]\label{assume:bk-model}
中介和结果的条件期望满足线性模型：
\[
\begin{cases}
\mathbb{E}(M \mid Z, X) = \beta_0 + \beta_1 Z + \beta_2^\top X, \\[4pt]
\mathbb{E}(Y \mid Z, M, X) = \theta_0 + \theta_1 Z + \theta_2 M + \theta_4^\top X.
\end{cases}
\tag{27.1}
\]
\end{assumption}

在上述线性模型下，NDE 和 NIE 可以简化为系数的简洁函数。

\begin{Corollary}[Baron--Kenny 公式]\label{cor:bk-formulas}
在 \cref{assume:mediation-no-z-y-confounding,assume:mediation-no-m-y-confounding,assume:mediation-no-z-m-confounding,assume:mediation-crossworld,assume:bk-model} 下，有
\[
\NDE = \theta_1, \qquad \NIE = \theta_2 \beta_1.
\tag{27.2}
\]
\end{Corollary}

\fboxsep=5pt\fboxrule=1pt
\noindent\fbox{
\begin{minipage}{0.97\textwidth}
\textbf{与 DAG 的对应}：在 \cref{fg:mediation-dag-bk} 中，$\theta_1$ 是路径 $Z \to Y$ 的系数（直接效应），$\beta_1 \theta_2$ 是路径 $Z \to M \to Y$ 的系数之积（间接效应）。这与 \cref{cor:bk-formulas} 完全一致。
\end{minipage}
}

\fboxsep=5pt\fboxrule=1pt
\noindent\fbox{
\begin{minipage}{0.97\textwidth}
\textbf{为何不需要对 $X$ 求和？}注意在 \cref{eq:nde-formula} 和 \cref{eq:nie-formula} 中，$\NDE(x)$ 和 $\NIE(x)$ 是在给定 $X=x$ 的条件下的值。而在 Baron--Kenny 线性模型下，$\theta_1$ 和 $\theta_2 \beta_1$ 与 $x$ 无关，因此条件效应等于无条件效应。
\end{minipage}
}

Baron--Kenny 公式的推导见附录 \cref{sec:proof-bk-formulas}。

如果获得了系数的 OLS 估计量 $\hat\beta_1$ 和 $\hat\theta_2$，则 NDE 和 NIE 的估计量为
\[
\hat{\NDE} = \hat\theta_1, \qquad \hat{\NIE} = \hat\theta_2 \hat\beta_1,
\]
这就是经典的 Baron--Kenny 估计。\Citet{sobel1982asymptotic} 和 \Citet{sobel1986some} 利用 delta 方法给出了 $\hat{\NIE}$ 的标准误：
\[
\hat{\text{var}}(\hat{\NIE}) = \hat{\text{var}}(\hat\theta_2) \hat\beta_1^2 + \hat\theta_2^2 \hat{\text{var}}(\hat\beta_1).
\]
基于此的假设检验称为 Sobel's test。\footnote{关于 Sobel 检验的更多讨论见 \cite{MacKinnon::2002}。}

\fboxsep=5pt\fboxrule=1pt
\noindent\fbox{
\begin{minipage}{0.97\textwidth}
\textbf{Baron--Kenny 方法的历史}：虽然 \Citet{baron1986moderator} 使这一方法广为人知，但它实际上有更早的渊源，包括 \Citet{hyman1955survey}、\Citet{alwin1975decomposition} 和 \Citet{Judd::1981} 的工作。\Citet{sobel1982asymptotic} 进一步发展了推断方法。
\end{minipage}
}

\section{应用实例：JOBS II 数据}\label{sec:bk-jobs}

我们通过 R 代码实现 Baron--Kenny 方法，并将其应用于 JOBS II 数据集。\footnote{JOBS II 是一个关于就业培训对失业工人影响的随机化现场实验。数据集在 R 的 \texttt{mediation} 包中。}

\begin{Example}[JOBS II：就业培训对心理健康的中介分析]\label{eg:jobs-mediation}
我们关注就业培训（处理 $Z$）对心理健康（结果 $Y$，用抑郁量表得分度量）的影响，以及这一影响是否通过工作搜寻效能（中介 $M$）传导。背景协变量 $X$ 包括经济困难程度、基线抑郁得分、性别、年龄、职业、婚姻状态、种族、教育程度和家庭收入。

通过以下 R 代码实现 Baron--Kenny 方法：

\begin{verbatim}
library("car")
BKmediation = function(Z, M, Y, X) {
  mediator.reg   = lm(M ~ Z + X)
  mediator.Zcoef = mediator.reg$coef[2]
  mediator.Zse   = sqrt(hccm(mediator.reg)[2, 2])

  outcome.reg    = lm(Y ~ Z + M + X)
  outcome.Zcoef  = outcome.reg$coef[2]
  outcome.Zse    = sqrt(hccm(outcome.reg)[2, 2])
  outcome.Mcoef  = outcome.reg$coef[3]
  outcome.Mse    = sqrt(hccm(outcome.reg)[3, 3])

  NDE = outcome.Zcoef
  NIE = outcome.Mcoef * mediator.Zcoef

  NDE.se = outcome.Zse
  NIE.se = sqrt(outcome.Mse^2 * mediator.Zcoef^2 +
                  outcome.Mcoef^2 * mediator.Zse^2)

  res = matrix(c(NDE, NIE, NDE.se, NIE.se,
                 NDE/NDE.se, NIE/NIE.se), 2, 3)
  rownames(res) = c("NDE", "NIE")
  colnames(res) = c("est", "se", "t")
  res
}
\end{verbatim}

运行结果如下：
\begin{verbatim}
> res = BKmediation(Z, M, Y, X)
> round(res, 3)
       est      se       t
NDE -0.037  0.042  -0.885
NIE -0.014  0.009  -1.528
\end{verbatim}

估计的直接效应 $\hat{\NDE} = -0.037$（培训对心理健康的直接保护效应）和间接效应 $\hat{\NIE} = -0.014$（通过提高工作搜寻效能而改善心理健康）均为负，但统计上均不显著。这表明 JOBS II 培训项目对心理健康的效应尚无充分证据支持按中介路径分解。
\end{Example}

\fboxsep=5pt\fboxrule=1pt
\noindent\fbox{
\begin{minipage}{0.97\textwidth}
\textbf{一个值得注意的现象}：在某些设定下，间接效应可能出现违反直觉的符号。例如，如果 $Z \to M$ 的效应 $\beta_1 > 0$，$M \to Y$ 的效应 $\theta_2 > 0$，但 $Z$ 与 $M$ 之间存在负交互项 $\theta_3 < 0$，则乘积 $\beta_1 (\theta_2 + \theta_3)$ 可能为负。这被称为\textbf{中介悖论}（mediator paradox）或\textbf{替代终点悖论}（surrogate endpoint paradox）。\footnote{详见 \cite{chen2007criteria}。}
\end{minipage}
}

\section{敏感性分析}

中介分析依赖于强且不可检验的假设。关键的不可检验假设是跨世界独立性（\cref{assume:mediation-crossworld}）：我们无法通过任何实验数据验证 $Y(z, m)$ 与 $M(z')$ 在 $z \neq z'$ 时是否独立。

因此，研究者常常需要进行\textbf{敏感性分析}（sensitivity analysis），以探索当关键假设不成立时，结论会如何变化。主要文献包括：

\begin{itemize}
  \item \Citet{ding2016sharp} 提出了类似 Cornfield 的敏感性界；
  \item \Citet{zhang2022interpretable} 提出了针对 Baron--Kenny 线性模型量身定制的敏感性分析方法。
\end{itemize}

敏感性分析的具体方法超出了本章的范围，但读者应记住：中介分析的结果必须谨慎解读，因为其核心假设无法被数据验证。

\section{本章小结}

本章围绕中介分析展开，主要内容如下：

\begin{enumerate}
  \item \textbf{核心问题}：因果效应如何在直接通路和间接通路之间分配？
  \item \textbf{嵌套潜在结果} $Y(z, M_{z'})$：描述跨世界 counterfactual，是定义 NDE 和 NIE 的基础。
  \item \textbf{效应分解}：在 Composition 假设下，总效应 $\tau = \NDE + \NIE$。
  \item \textbf{识别假设}：四个假设（序列可忽略性 + 跨世界独立性）下，中介公式成立。
  \item \textbf{中介公式}：将不可观测的 NDE 和 NIE 表达为可观测条件分布的函数。
  \item \textbf{Baron--Kenny 公式}：在 \cref{assume:bk-model} 下，$\NDE = \theta_1$，$\NIE = \theta_2 \beta_1$。
  \item \textbf{跨世界 counterfactual 的哲学争议}：严格的 Popperian 认识论者认为中介分析是形而上学的。
  \item \textbf{敏感性分析}：由于核心假设不可检验，敏感性分析是必要的。
\end{enumerate}

\par\Notes{中介公式的完整证明见附录 \cref{sec:proof-mediation}；Baron--Kenny 公式的推导见附录 \cref{sec:proof-bk-formulas}；各节练习题见附录 \cref{sec:mediation-exercises}。}

\section*{练习}\label{sec:mediation-exercises}

% === 练习内容 ===

\begin{Exercise}\label{ex:mediation-robins-latent-groups}
%（\cite{Robins::1992} 的潜在群组澄清）

假设 $(Z, M, Y)$ 均为二值变量。基于潜在结果
\[
M(1), M(0), Y(1, M_1), Y(1, M_0), Y(0, M_1), Y(0, M_0)
\]
的联合取值，单元可以分成多少类型？

若进一步假设单调性 $M(1) \geq M(0)$（$Z$ 对 $M$ 的单调性），则单元可以分成多少类型？

若再假设 $Y(1, m) \geq Y(0, m)$ 且 $Y(z, 1) \geq Y(z, 0)$ 对所有 $z, m$ 成立（$Z$ 和 $M$ 对 $Y$ 的单调性），则单元可以分成多少类型？
\end{Exercise}

\par\Notes{本题答案见附录 \cref{sec:sol-mediation-robins-latent-groups}。}

\begin{Exercise}\label{ex:sr-jr-equivalence}
%（序列随机化与联合随机化的等价性）

证明 \cref{eq:joint-randomization} 与 \cref{assume:mediation-no-z-y-confounding,assume:mediation-no-m-y-confounding} 等价。
\end{Exercise}

\par\Notes{本题答案见附录 \cref{sec:sol-sr-jr}。}

\begin{Exercise}\label{ex:alternative-si}
%（\cite{Imai::2010} 的另一组假设）

\Citet{Imai::2010} 提出了另一组假设来推导中介公式：
\[
\{Y(z, m), M(z')\} \ind Z \mid X, \qquad Y(z, m) \ind M(z') \mid (Z = z', X),
\]
对所有 $z, z', m$ 成立。

证明：在上述假设下，\cref{eq:mediation-formula} 中的中介公式同样成立。
\end{Exercise}

\par\Notes{本题答案见附录 \cref{sec:sol-alternative-si}。}

\begin{Exercise}\label{ex:difference=product}
%（差值法与乘积法的等价性）

在 Baron--Kenny 线性模型（\cref{assume:bk-model}）下，首先拟合：
\[
\hat M = \hat\beta_0 + \hat\beta_1 Z + \hat\beta_2^\top X, \qquad
\hat Y = \hat\theta_0 + \hat\theta_1 Z + \hat\theta_2 M + \hat\theta_4^\top X.
\]
则 NIE 的估计量为乘积 $\hat\theta_2 \hat\beta_1$（乘积法）。

另一方面，首先拟合：
\[
\hat Y = \hat\alpha_0 + \hat\alpha_1 Z + \hat\alpha_1^\top X,
\]
再拟合：
\[
\hat Y = \tilde\theta_0 + \tilde\theta_1 Z + \tilde\theta_2 M + \tilde\theta_4^\top X,
\]
则 NIE 的另一估计量为差值 $\hat\alpha_1 - \tilde\theta_1$（差值法）。

证明：$\hat\alpha_1 - \tilde\theta_1 = \hat\theta_2 \hat\beta_1$。
\end{Exercise}

\par\Notes{本题答案见附录 \cref{sec:sol-difference-product}。}

\begin{Exercise}\label{ex:nie-binary-m}
%（二值中介的 NIE）

证明 \cref{cor:nie-binary}（二值中介的 NIE 乘积公式）。
\end{Exercise}

\par\Notes{本题答案见附录 \cref{sec:sol-nie-binary}。}

\begin{Exercise}\label{ex:interaction}
%（处理—中介交互项）

\Citet{vanderweele2015explanation} 建议在结果模型中加入处理与中介的交互项：
\[
\begin{cases}
\mathbb{E}(M \mid Z, X) = \beta_0 + \beta_1 Z + \beta_2^\top X, \\[4pt]
\mathbb{E}(Y \mid Z, M, X) = \theta_0 + \theta_1 Z + \theta_2 M + \theta_3 ZM + \theta_4^\top X.
\end{cases}
\]

证明：在此模型下，
\[
\NDE = \theta_1 + \theta_3\{\beta_0 + \beta_2^\top \mathbb{E}(X)\}, \qquad
\NIE = (\theta_2 + \theta_3) \beta_1.
\]
并说明如何在 IID 数据下估计 $\NDE$ 和 $\NIE$。
\end{Exercise}

\par\Notes{本题答案见附录 \cref{sec:sol-interaction}。}

\begin{Exercise}\label{ex:contin-m-binary-y}
%（连续中介 + 二值结果）

考虑中介 $M$ 服从正态线性模型、结果 $Y$ 服从 logistic 模型：
\[
\begin{cases}
M \mid Z, X \sim \mathcal{N}(\beta_0 + \beta_1 Z + \beta_2^\top X, \sigma_M^2), \\[4pt]
\operatorname{logit}\{\Pr(Y = 1 \mid Z, M, X)\} = \theta_0 + \theta_1 Z + \theta_2 M + \theta_4^\top X.
\end{cases}
\]
请将 $\NDE$ 和 $\NIE$ 表示为模型参数和 $X$ 分布的函数，并说明如何在 IID 数据下估计它们。
\end{Exercise}

\par\Notes{本题答案见附录 \cref{sec:sol-contin-m-binary-y}。}

\begin{Exercise}\label{ex:logit-binary-m}
%（二值中介 + 连续结果）

考虑二值中介 $M$ 服从 logistic 模型、结果 $Y$ 服从线性模型：
\[
\begin{cases}
\operatorname{logit}\{\Pr(M = 1 \mid Z, X)\} = \beta_0 + \beta_1 Z + \beta_2^\top X, \\[4pt]
\mathbb{E}(Y \mid Z, M, X) = \theta_0 + \theta_1 Z + \theta_2 M + \theta_4^\top X.
\end{cases}
\]

证明：
\[
\NDE = \theta_1, \qquad
\NIE = \theta_2 \, \mathbb{E}\big\{ \operatorname{expit}(\beta_0 + \beta_1 + \beta_2^\top X) - \operatorname{expit}(\beta_0 + \beta_2^\top X) \big\}.
\]
并说明如何在 IID 数据下估计 $\NDE$ 和 $\NIE$。
\end{Exercise}

\par\Notes{本题答案见附录 \cref{sec:sol-logit-binary-m}。}

\begin{Exercise}\label{ex:no-crossworld-modification}
%（放松跨世界独立性假设）

将嵌套潜在结果修改为
\[
Y(z, F_{M_{z'} \mid X}) = \int Y(z, m) \, f(M_{z'} = m \mid X) \, \mathrm{d}m,
\]
其中 $F_{M_{z'} \mid X}$ 是条件分布而非确定值。相应地将 NDE 和 NIE 定义为
\[
\NDE = \mathbb{E}\{ Y(1, F_{M_0 \mid X}) - Y(0, F_{M_0 \mid X}) \}, \quad
\NIE = \mathbb{E}\{ Y(1, F_{M_1 \mid X}) - Y(1, F_{M_0 \mid X}) \}.
\]

证明：在仅假设 \cref{assume:mediation-no-z-y-confounding,assume:mediation-no-m-y-confounding,assume:mediation-no-z-m-confounding}（不要求跨世界独立性）下，$\NDE$ 和 $\NIE$ 的识别公式与正文相同。
\end{Exercise}

\par\Notes{本题答案见附录 \cref{sec:sol-no-crossworld}。}

\begin{Exercise}\label{ex:principal-mediation-connection}
%（主分层与中介分析的联系）

\Citet{vanderweele2008simple} 和 \Citet{forastiere2018principal} 综述并比较了主分层与中介分析两种方法。请查阅这些文献，并撰写一页总结，阐述两种方法的异同与各自的优缺点。
\end{Exercise}

\par\Notes{本题为开放性研究型练习，请参考 \cite{vanderweele2008simple} 和 \cite{forastiere2018principal}。}

\endinput
